\section{Related Work}
\label{sec:related_work}

Since the publication of the Transformer architecture \cite{vaswani2017attention} and BERT \cite{devlin2018bert}, the field of NLP has experienced a massive shift towards the use of LLMs with self-supervised pre-training. 
Multiple masked langauge models, including T5 \cite{raffel2020exploring} and Megatron-LM \cite{shoeybi2019megatron}, have shown consistent improvements through scale. These scaling gains come not only from growing the total number of parameters in the models, but also the amount and quality of pre-training data~\cite{liu2019roberta,chinchilla2022}.  

Auto-regressive language models~\cite{mikolov2009neural} have seen the largest growth in model size, from 117M parameters \cite{radford2018gpt} to over 500B parameters \cite{megatron2022,palm2022}. The resulting massive improvement in generative fluency and quality was first characterized in GPT-2 \cite{radford2019language} and further improved with \biggpt{} \cite{brown2020gpt3} and later models. Although a variety of very large (over 100B parameters) generative models have now been trained \cite{J1WhitePaper,gopher2022,thoppilan2022lamda,megatron2022,palm2022}, they are all closed source and accessible only internally or via paid API services. There are a few notable efforts towards open sourcing LLMs from non-profit research organizations including EleutherAI~\cite{eleutherai2022} and BigScience.\footnote{\url{https://huggingface.co/bigscience/tr11-176B-ml-logs/tensorboard}} These models differ from the OPT models in pre-training data, target languages and model scale, making it possible for the community to compare different pre-training strategies.


Since \citet{brown2020gpt3}, the primary evaluation criterion for LLMs has been prompt-based \cite{eleutherai2022,gopher2022,palm2022}, as is also performed in this paper. This is largely due to the convenience of evaluating on many tasks without specialized task-specific fine-tuning. Prompting itself has a long history: cloze evaluations go back several decades \cite{chambers-jurafsky-2008-unsupervised,storycloze}. More recently, prompting or masked infilling has been used to probe models for knowledge \cite{petroni2019language} or perform a variety of NLP tasks \cite{radford2019language, brown2020gpt3}.
There has also been work on eliciting prompting behavior in smaller models \cite{timopriming2020,gao2021,Li2021Prefix-Tuning:Generation,DBLP:journals/corr/abs-2104-08691,Scao2021HowWorth}, improving the flexibility of prompting \cite{Shin2020AutoPrompt}, and understanding why and how prompting works \cite{incontextgpt3_2021,min2022rethinking}. 

Recent efforts have shown gains by fine-tuning models to directly respond to instruction-style prompting \cite{flan2021,min2021metaicl,sanh2021multitask,ouyang2022training}. However, effective prompt engineering remains an open research challenge. Results vary significantly and unpredictably with the selection of the prompt  \cite{promptsensitivity2021}, and models do not seem to understand the prompts as fully as we expect \cite{webson2021prompt}. Furthermore, it is challenging to write prompts without a development set, which leads to questions about the extent to which we are actually achieving zero- or few-shot learning in practice~\cite{perez2021true}. We do not attempt to address these concerns of prompting, and instead only aim to provide evaluation of \OPT{} in existing settings. However, we hope the full release of \OPT{} will enable others to better study these challenges in the future. 